\section{The classic arguments}\label{sec:arguments}

This section applies T-Class (\cref{item:TClass}) case by case. For each argument it states a
standard form, names the \emph{decisive element}, the one that carries the conclusion beyond the
evidence, and gives that element's class: an account, a judgment, or an imposed constraint. By
P-M(b), a decisive premise offered \emph{a priori}, in an argument valid for its conclusion, would
exclude an account consistent with the truths of reason, so by T-Class it is an imposed constraint.
In a non-deductive argument, the rule of inference is a judgment.

\subsection{A worked example}\label{ss:worked-example}

One argument is set out here in full, to show how the classification proceeds; the others are given
more briefly. The conclusions this example previews are derived in \cref{sec:result}.

The fine-tuning argument, in its probabilistic form, has four elements
\citep{swinburne2004,collins2009}. First, the constants of nature, such as the strength of gravity
and the cosmological constant, fall in a range that permits life, and on the argument's estimates a
narrow one. Second, a life-permitting universe is far more probable if it was designed than on its
rivals: on design it is to be expected, and on a rival on which the constants have no explanation it
is a coincidence. Third, design has a prior probability that is not too low. Fourth, by Bayes'
theorem, the probability of an account given the evidence rises with how much more probable the
evidence is on that account than on its rivals, and with the account's prior probability; so, given
the first three elements, design is probable.

Each element falls in one of T-Class's classes. The first is evidence: it reports the values of the
constants as physics has measured them. The second is a judgment: it passes from the evidence to a
comparison of accounts by estimating how probable the evidence is on each, and probabilistic
confirmation is a non-deductive rule. The third is an imposed constraint: a prior ranks the accounts
before the evidence is weighed, and it is a criterion for choosing among them that the evidence does
not determine. The fourth is a valid inference, since Bayes' theorem is a theorem of probability, and
so it is a truth of reason.

Design and its rivals do not dispute the evidence, and the inference is valid, so neither of these
carries the conclusion beyond the evidence. The decisive elements are the judgment over likelihoods
and the prior.

Each is answered by a different conclusion. Against a claim that the argument establishes design,
C1 applies: an account with no designer fits the same evidence, and design is false in it. One such
account has many universes with varying constants, with observers found only in those that permit
life; another derives the constants from a deeper law. The argument as stated concludes only that
design is probable, and that conclusion C1 leaves standing. That conclusion rests on the prior as well
as the likelihoods, and for the prior C3 applies. A prior is a choice the evidence does not determine. A prior low enough, equally
consistent with the evidence, makes design improbable on the same likelihoods; one high enough makes
it probable.

C2 applies in one further case. If design and its rival agree on all possible evidence, so that
every observation that could ever be made is as probable on one as on the other, then no further
measurement separates them. If they differ somewhere on possible evidence, the case is C1's.

\subsection{Arguments for theism}\label{ss:arguments-for}

\paragraph{The kalam argument.} Whatever begins to exist has a cause; the universe began to exist;
so the universe has a cause \citep{craig1979}. Craig supports the second premise with the claim
that an actual infinite cannot exist, and argues in a further step that the cause is personal. The
decisive elements are the causal principle, the denial of an actual infinite, and the step to a
personal cause. Each is an account if offered as a generalization from experience, and an imposed
constraint if offered \emph{a priori}.

\paragraph{The contingency argument.} Every contingent fact has an explanation; the totality of
contingent facts cannot be explained by a contingent being; so there is a necessary being
\citep{pruss2009}. The decisive elements are the principle of sufficient reason and the claim that
the totality has no contingent explanation, classed as for the kalam argument.

\paragraph{The Nyāya argument.} Udayana argues that the earth, being an effect, has an intelligent
maker, Īśvara, as a pot has a potter \citep{dasti2011}. The decisive element is that the earth is an
effect requiring an intelligent maker, an account if drawn from experience and an imposed
constraint if offered \emph{a priori}.

\paragraph{The fine-tuning argument.} Its probabilistic form is set out and classified in
\cref{ss:worked-example}.

\paragraph{The moral argument.} There are objective moral facts, and a theory on which they consist
in facts about God, such as God's commands, explains them better than its nontheistic rivals
\citep{adams1987}. The decisive element is that objective moral facts exist and need a ground. It is
an account if offered as a description of how things are, and an imposed constraint if offered
\emph{a priori}.

\paragraph{The ontological argument.} In its modal form: it is possible that a maximally great being
exists; maximal greatness includes necessary existence; in the modal system S5, whatever is possibly
necessary is necessary; so a maximally great being exists \citep{plantinga1974}; compare
\citet[ch.~2--3]{anselm-proslogion}. The decisive elements are the possibility premise and the modal
axiom S5. Both are offered \emph{a priori}, and so are imposed constraints.

\paragraph{The arguments from consciousness and reason; the trademark argument.} Mind, rational
insight, or the idea of an infinite being cannot arise from unthinking nature, so they come from a
mind; Descartes argues that the idea of the infinite must have been placed in us by an infinite cause
\citep[Meditation~III]{descartes1641}. The decisive element is that mind, or the idea of the
infinite, cannot arise from nature: an account, or an imposed constraint if offered as a causal
principle held \emph{a priori}.

\paragraph{The argument from religious experience.} It seems to many people that God is present to
them, and by the principle of credulity what seems present probably is present, unless there is
reason to think otherwise \citep{swinburne2004}. Alston argues that it is practically rational to
engage in a socially established doxastic practice, the practice of forming beliefs on such
experiences no less than sense perception, and concedes that its reliability cannot be shown
without epistemic circularity \citep{alston1991}. The decisive elements are the principle of
credulity, an imposed constraint, and, for Alston, this parity principle of practical rationality,
a practical reason and so an imposed constraint.

\paragraph{Testimony and revelation.} Scripture and tradition transmit what God revealed, and the
transmission is faithful (for example, the Bible and the Qur'an; \citealp{bible,quran}). The
decisive element is that the transmission is reliable and its origin divine, an account.

\paragraph{Miracles, and Hume's argument against them.} For: God's action would raise the
probability of a reported event that nature does not explain \citep{swinburne2004}. Against: no
testimony suffices to establish a miracle unless its falsehood would be more miraculous than the
fact it reports \citep[§10.13]{hume1748}. The decisive element for the miracle is that God's action
raises the probability of the report, a judgment. The decisive element for Hume is his maxim, a
standard for preferring one conclusion to another, and so an imposed constraint.

\paragraph{First-person revelation.} The witness has the experience directly: Moses at the bush that
burned and was not consumed (Exodus 3:2--14), or Joseph Smith seeing two Personages in a pillar of
light \citep[1:16--17]{jsh1838}. The decisive element is that the experience is veridical, that it
was God who spoke, an account.

\paragraph{Inner witness and proper basicality.} A sense of the divine, or the inner testimony of the
Spirit, gives rise to belief in God directly, without inference \citep[I.3.1, I.7.4]{calvin1559},
and a belief so formed, by a faculty working as designed, may be properly basic
\citep{plantinga2000}. The decisive element is that the faculty is reliable, an account.

\paragraph{\emph{Consensus gentium}.} Belief in God is held across peoples and ages, and such
widespread agreement is offered as evidence of its truth; \citet{kelly2011} examines a modest version
of the argument. The decisive element is that convergence of belief tracks the truth, an account.

\paragraph{Pascal's Wager; James's will to believe.} Where the evidence leaves the question open, the
expected value of believing favors belief \citep{pascal1670}; or, where the option is live, forced
and momentous, one has the right to let one's passional nature decide \citep{james1896}. The
decisive element is expected value, or the right to believe, an imposed constraint (a choice).

\subsection{Arguments against theism}\label{ss:arguments-against}

\paragraph{The logical problem of evil.} A good thing eliminates evil as far as it can, and there are
no limits to what an omnipotent being can do; so an omnipotent, wholly good being would eliminate
evil completely; evil exists; so there is no such being \citep{mackie1955}. The decisive element is
that God and evil are incompatible, offered \emph{a priori}, and so an imposed constraint.

\paragraph{The evidential problem of evil.} There exist instances of intense suffering that an
omnipotent, omniscient being could have prevented without losing some greater good or permitting
some evil equally bad or worse, and such a being, if wholly good, would prevent them
\citep{rowe1979}. In a second form, the facts about pain and pleasure are more probable on the
hypothesis of indifference than on theism \citep{draper1989}. The decisive element for Rowe is that
such suffering exists, an account. For Draper it is the comparison of probabilities, a judgment,
with the prior an imposed constraint.

\paragraph{Divine hiddenness.} A perfectly loving God would ensure that everyone capable of a
relationship with him is in a position to believe that he exists, so there would be no reasonable,
or inculpable, nonbelief; there is such nonbelief; so there is no such God \citep{schellenberg1993}.
The decisive element is that perfect love entails that there is no reasonable nonbelief. It is an
account if it states what love is in fact like, and an imposed constraint if offered
\emph{a priori}.

\paragraph{Findlay's disproof.} An adequate object of worship must exist necessarily, and no
existential proposition can be necessary, so there is no such object \citep{findlay1948}. The
decisive element is that no existential proposition can be necessary, offered \emph{a priori}, and
so an imposed constraint.

\paragraph{Occam's razor.} An account of the world without God explains the same evidence with fewer
entities, so it is to be preferred (on such appeals to parsimony, see \citealp{sober2015}). The
decisive element is that the simpler account is to be preferred, a standard for preferring one
conclusion to another, and so an imposed constraint.

\subsection{Naturalism and the skeptical hypothesis}\label{ss:arguments-naturalism}

\paragraph{Causal closure.} Every physical effect has a sufficient physical cause, and Papineau
argues from this to physicalism \citep{papineau2001}; extended, it leaves nothing beyond nature that
acts in the physical world. The decisive element is that every physical effect has a sufficient
physical cause, an account.

\paragraph{The evolutionary argument against naturalism.} If naturalism and evolution are both true,
the probability that our cognitive faculties are reliable is low or inscrutable; so whoever accepts
both has a defeater for every belief, including naturalism \citep[ch.~12]{plantinga1993}. The
decisive elements are that this probability is low or inscrutable, an account, and the principle
that such a probability defeats the beliefs the faculties produce, an imposed constraint.

\paragraph{The simulation argument.} Either almost no civilizations reach the stage at which they can
run detailed simulations of minds like ours, or almost none that reach it choose to, or we are almost
certainly living in a simulation \citep{bostrom2003}. The decisive elements are the assumptions
about simulators, that minds like ours can be simulated and what mature civilizations could and
would do, an account; and the bland indifference principle, that one's credence in being simulated
should match the fraction of observers like oneself who are, an imposed constraint.
